% GENERATED FILE. Edit canonical lessons, then run npm run generate:books.
% canonical-source-hash: fnv1a64:84313318f0e39024
% canonical-lessons: KA-C299-muduki, KA-C299-jana, KA-C299-manusya, KA-C299-aparicita, KA-C299-hiriyaru

\chapter{Old woman, People, Human being, Stranger, Elders}
\label{ch:ka-old-woman-people-human-being-stranger-elders}
\hlchaptermodality{299}

\begin{chapteropening}
\textbf{By the end of this chapter:} \emph{I can say an old woman, people, a human being, a stranger and elders.}

Complete the last lesson of chapter 299: I can say an old woman, people, a human being, a stranger and elders.
\end{chapteropening}

\section[\texorpdfstring{\kn{ಮುದುಕಿ} (muduki)}{ಮುದುಕಿ (muduki)}]{\kn{ಮುದುಕಿ} (muduki) — an old woman}

\label{lesson:KA-C299-muduki}



\begin{quote}
Before the new one: say the Kannada for a competition, then the Kannada for a prize.
\end{quote}

\subsection*{You'll want to know: \kn{ಮುದುಕಿ}}
\textbf{\kn{ಮುದುಕಿ}} — \emph{muduki} — \textquotedblleft{}an old woman\textquotedblright{}.

\begin{grammarlens}[title={What you've built}]
One more word for this chapter.
\end{grammarlens}

\subsection*{Your turn}
\begin{itemize}
\raggedright
  \item \emph{Say it:} \emph{muduki}
  \item \emph{Say it:} \emph{muduki}, once more
  \item \emph{Say it:} say \emph{muduki}
  \item \emph{Recall:} say the Kannada for a competition, then the Kannada for a prize, then say \emph{muduki} again
\end{itemize}

\begin{tcolorbox}[breakable,colback=teal!4,colframe=teal!35!black,title={Before you move on}]
What does \kn{ಮುದುಕಿ} mean? (\textquotedblleft{}An old woman\textquotedblright{}.) Say it once more.
\end{tcolorbox}
\section[\texorpdfstring{\kn{ಜನ} (jana)}{ಜನ (jana)}]{\kn{ಜನ} (jana) — people}

\label{lesson:KA-C299-jana}



\begin{quote}
Before the new one: say the Kannada for a prize, then the Kannada for an old woman.
\end{quote}

\subsection*{You'll want to know: \kn{ಜನ}}
\textbf{\kn{ಜನ}} — \emph{jana} — \textquotedblleft{}people\textquotedblright{}.

\begin{grammarlens}[title={What you've built}]
One more word for this chapter.
\end{grammarlens}

\subsection*{Your turn}
\begin{itemize}
\raggedright
  \item \emph{Say it:} \emph{jana}
  \item \emph{Say it:} \emph{jana}, once more
  \item \emph{Say it:} say \emph{jana}
  \item \emph{Recall:} say the Kannada for a prize, then the Kannada for an old woman, then say \emph{jana} again
\end{itemize}

\begin{tcolorbox}[breakable,colback=teal!4,colframe=teal!35!black,title={Before you move on}]
What does \kn{ಜನ} mean? (\textquotedblleft{}People\textquotedblright{}.) Say it once more.
\end{tcolorbox}
\section[\texorpdfstring{\kn{ಮನುಷ್ಯ} (manuṣya)}{ಮನುಷ್ಯ (manuṣya)}]{\kn{ಮನುಷ್ಯ} (manuṣya) — a human being}

\label{lesson:KA-C299-manusya}



\begin{quote}
Before the new one: say the Kannada for an old woman, then the Kannada for people.
\end{quote}

\subsection*{You'll want to know: \kn{ಮನುಷ್ಯ}}
\textbf{\kn{ಮನುಷ್ಯ}} — \emph{manuṣya} — \textquotedblleft{}a human being\textquotedblright{}.

\begin{grammarlens}[title={What you've built}]
One more word for this chapter.
\end{grammarlens}

\subsection*{Your turn}
\begin{itemize}
\raggedright
  \item \emph{Say it:} \emph{manuṣya}
  \item \emph{Say it:} \emph{manuṣya}, once more
  \item \emph{Say it:} say \emph{manuṣya}
  \item \emph{Recall:} say the Kannada for an old woman, then the Kannada for people, then say \emph{manuṣya} again
\end{itemize}

\begin{tcolorbox}[breakable,colback=teal!4,colframe=teal!35!black,title={Before you move on}]
What does \kn{ಮನುಷ್ಯ} mean? (\textquotedblleft{}A human being\textquotedblright{}.) Say it once more.
\end{tcolorbox}
\section[\texorpdfstring{\kn{ಅಪರಿಚಿತ} (aparicita)}{ಅಪರಿಚಿತ (aparicita)}]{\kn{ಅಪರಿಚಿತ} (aparicita) — a stranger}

\label{lesson:KA-C299-aparicita}



\begin{quote}
Before the new one: say the Kannada for people, then the Kannada for a human being.
\end{quote}

\subsection*{You'll want to know: \kn{ಅಪರಿಚಿತ}}
\textbf{\kn{ಅಪರಿಚಿತ}} — \emph{aparicita} — \textquotedblleft{}a stranger\textquotedblright{}.

\begin{grammarlens}[title={What you've built}]
One more word for this chapter.
\end{grammarlens}

\subsection*{Your turn}
\begin{itemize}
\raggedright
  \item \emph{Say it:} \emph{aparicita}
  \item \emph{Say it:} \emph{aparicita}, once more
  \item \emph{Say it:} say \emph{aparicita}
  \item \emph{Recall:} say the Kannada for people, then the Kannada for a human being, then say \emph{aparicita} again
\end{itemize}

\begin{tcolorbox}[breakable,colback=teal!4,colframe=teal!35!black,title={Before you move on}]
What does \kn{ಅಪರಿಚಿತ} mean? (\textquotedblleft{}A stranger\textquotedblright{}.) Say it once more.
\end{tcolorbox}
\section[\texorpdfstring{\kn{ಹಿರಿಯರು} (hiriyaru)}{ಹಿರಿಯರು (hiriyaru)}]{\kn{ಹಿರಿಯರು} (hiriyaru) — elders}

\label{lesson:KA-C299-hiriyaru}



\begin{quote}
Before the new one: say the Kannada for a human being, then the Kannada for a stranger.
\end{quote}

\subsection*{You'll want to know: \kn{ಹಿರಿಯರು}}
\textbf{\kn{ಹಿರಿಯರು}} — \emph{hiriyaru} — \textquotedblleft{}elders\textquotedblright{}.

\begin{grammarlens}[title={What you've built}]
One more word for this chapter.
\end{grammarlens}

\subsection*{Your turn}
\begin{itemize}
\raggedright
  \item \emph{Say it:} \emph{hiriyaru}
  \item \emph{Say it:} \emph{hiriyaru}, once more
  \item \emph{Say it:} say \emph{hiriyaru}
  \item \emph{Recall:} say the Kannada for a human being, then the Kannada for a stranger, then say \emph{hiriyaru} again
\end{itemize}

\begin{tcolorbox}[breakable,colback=teal!4,colframe=teal!35!black,title={Before you move on}]
What does \kn{ಹಿರಿಯರು} mean? (\textquotedblleft{}Elders\textquotedblright{}.) Say it once more.
\end{tcolorbox}
